\documentclass{article}
\usepackage[final]{neurips_2026}

\usepackage[utf8]{inputenc}
\usepackage[T1]{fontenc}
\usepackage{hyperref}
\usepackage{url}
\usepackage{booktabs}
\usepackage{amsfonts}
\usepackage{microtype}
\usepackage{xcolor}
\usepackage{graphicx}
\usepackage{amsmath}

% TODO markers for content not yet finalized (figures pending scripts/, citations pending a lit pass).
\newcommand{\todo}[1]{\textcolor{red}{[TODO: #1]}}

\title{The Execution Floor: Budget Allocation, Not Scaffolding, Search, or Training,
Moves Whole-Proof Theorem Provers}

\author{%
  \todo{author list, anonymize for submission} \\
}

\begin{document}

\maketitle

\begin{abstract}
For a fixed whole-proof theorem prover, we ask whether test-time interventions --- agentic scaffolding
(premise retrieval, failure memory, an LLM reviewer, tactic-skeleton hints, forced approach diversity),
additional search, or additional per-problem budget --- can move solve rate beyond what raw sampling
already buys at a given token budget $B$. Across two independently-trained provers (Goedel-Prover-V2-8B,
DeepSeek-Prover-V2-7B), two benchmarks (in-distribution miniF2F, out-of-distribution ProofNet\#), and
eight experimental phases, we find a consistent negative result: none of the above move solve rate. A
pre-registered, causally interventional test (forcing approach diversity on the fully-trapped problem
core) confirms the mechanism directly --- diversity rises 42--70\% under a budget-matched manipulation
check, yet solves stay flat and output quality mildly degrades. We extend the null to \emph{training}:
neither closing-likelihood SFT, lightweight post-hoc RL, nor full-pipeline lab-scale RL (two independently
matched Base$\to$SFT$\to$RL lineages) moves the floor either. The one lever that reliably moves outcomes
at fixed compute is \emph{how} a budget is allocated across a batch of \emph{different} problems, not
any change to any single problem's search or training: a realizable, mechanism-informed allocation policy
saves \todo{XX}\% of compute at 90\% of full-budget accuracy on one model (Goedel$\times$ProofNet\#,
per-seed robust), though this positive result is not yet confirmed on a second model. We position this
as a definitive-negative result against an active literature that treats scaffolding and test-time
compute as reliable levers, and as a methodological case study in verifier auditing, symmetric
skepticism of both surprising positives and surprising nulls, and causal (not just correlational)
mechanism testing.
\end{abstract}

\section{Introduction}
\label{sec:intro}

Large language models fine-tuned for formal theorem proving (``whole-proof'' provers that emit a
complete tactic proof in one generation, as opposed to step-by-step interactive search) have improved
rapidly, and a growing body of work treats \emph{test-time compute} --- more sampling, agentic
scaffolding, retrieval-augmented context, search over partial proof states --- as a reliable lever for
squeezing more solves out of a fixed model. We ask a narrower, falsifiable question: for two
independently-trained whole-proof provers, at a fixed per-problem token budget $B$, what actually moves
solve rate?

We find an \emph{execution floor}: a ceiling on solve rate that eight different interventions, spanning
search-time scaffolding, closing-tactic libraries, additional budget on already-attempted problems, and
three distinct training interventions (SFT, lightweight RL, full-pipeline RL), all fail to move. The one
exception is not a model or scaffold change at all, but a \emph{policy} question: given a fixed total
compute budget, how should it be split across a batch of different problems. This paper reports the
floor as the headline finding, with the allocation result as a constructive counterpoint scoped to what
is currently confirmed (Section~\ref{sec:allocation}).

\textbf{Contributions.}
\begin{itemize}
  \item A cross-model (2 provers), cross-benchmark (in-distribution / OOD) replication of the
    budget$\to$solve relationship and its saturation asymmetry (Section~\ref{sec:passb}).
  \item A one-factor-at-a-time ablation of six agentic scaffolding components showing none is a net
    positive lever, and that two are actively harmful out-of-distribution (Section~\ref{sec:scaffolding}).
  \item A four-part mechanism account (diversity collapse onto a deep-reasoning execution floor, on
    which the one retrieval variant we tested empirically fails) culminating in a \emph{pre-registered,
    causally interventional} test that confirms the mechanism directly rather than by correlation alone
    (Section~\ref{sec:mechanism}).
  \item An exhaustion sweep across five further candidate levers --- hammer/SMT tactics, more budget on
    already-trapped problems, closing-likelihood SFT, lightweight post-hoc RL, and full-pipeline
    lab-scale RL across two matched training lineages --- all null, extending the floor from
    search-time to training-time interventions (Section~\ref{sec:exhaustion}).
  \item A constructive counterpoint: compute-optimal budget allocation across problems, the one lever
    that does move outcomes, reported honestly as one-model-robust (Section~\ref{sec:allocation}).
  \item A methodology contribution: a verifier-soundness audit protocol, and the discipline of applying
    the same base-rate skepticism to surprising nulls as to surprising positives, which caught a
    verifier bug that would otherwise have produced a spurious ``training moves the floor'' headline
    (Section~\ref{sec:methodology}).
\end{itemize}

\section{Related work}
\label{sec:related}

\todo{Position against: (1) the test-time-compute scaling literature (more samples / more search buys
accuracy) --- we show this does not transfer to agentic scaffolding for whole-proof theorem proving;
(2) agentic/tool-augmented LLM provers and retrieval-augmented provers (ReProver and successors) ---
our ablation only tests BM25 retrieval, which fails on this failure distribution (F2,
Section~\ref{sec:mechanism}); neural (ReProver-style) retrieval remains untested here and we say so
explicitly, rather than let the BM25 result stand in for retrieval broadly; (3) step-by-step /
search-based provers (BFS-Prover, HTPS,
InternLM2.5-StepProver) --- out of scope by design (prior work found search $<$ Pass@1, and our
Section~\ref{sec:mechanism} findings say the bottleneck is execution depth, not search breadth), we test whole-proof
sampling only; (4) RL-for-reasoning post-training literature --- our Phase 6/8 results are a negative
data point against RL-moves-everything narratives, specifically for execution-limited domains; (5) the
self-correction-negative-result lineage, as precedent for reporting a well-controlled null as a
contribution in its own right. Full citation pass pending.}

\section{Setup}
\label{sec:setup}

\textbf{Models.} Two independently-trained whole-proof provers: Goedel-Prover-V2-8B
(\texttt{Goedel-LM/Goedel-Prover-V2-8B}, Lean~v4.9.0-rc1 + mathlib4) and DeepSeek-Prover-V2-7B
(\texttt{deepseek-ai/DeepSeek-Prover-V2-7B}, Lean~v4.9.0 + standard mathlib4). Pinned model
revisions, Lean toolchains, and mathlib commits are given in \todo{repro/PINS.md}.

\textbf{Benchmarks.} miniF2F-test (244 audited problems; competition-style, in-distribution for both
provers) and ProofNet\# (186 test; undergraduate-level, out-of-distribution). DeepSeek-Prover-V2
inherits Goedel's protocol exactly --- identical canonical statement sets, budgets, and agent
configuration, with only the model and Lean pin changed --- so any cross-model difference is
attributable to the model, not the harness.

\textbf{Verifier.} Lean itself (REPL backend) is the sole source of truth for ``solved.'' We treat the
verifier as a first-class audited artifact, not an assumed-correct oracle; Section~\ref{sec:methodology}
details two rounds of verifier-soundness bugs found and fixed over the project's lifetime, and their
measured blast radius on every reported number in this paper.

\textbf{Budget and variance protocol.} Token budgets $B \in \{2\text{k}, 8\text{k}, 32\text{k},
128\text{k}\}$; every headline number is 3 seeds minimum, reported mean $\pm$ seed-std. A mean delta
under $\sim$1 baseline-$\sigma$ is treated as noise throughout (the ``noise bar''); component-vs-baseline
comparisons use paired flips (per-(problem, seed) solved/unsolved transitions), not mean deltas, since
at $n=3$ seeds mean deltas can hide gains-equal-losses churn that a flip table exposes directly
(Section~\ref{sec:methodology}, lesson 1).

\textbf{pass@\emph{B} is not pass@\emph{N}.} We deliberately report pass@\emph{budget} (solved within a
cumulative token ledger shared across propose and refine calls for one problem) rather than the
test-time-compute literature's usual pass@\emph{N} (solved within $N$ independent, unrefined samples).
The two are not interchangeable, and the gap is budget-dependent, not a fixed offset: at $B$=2k the
\emph{median} cell across all four (model, benchmark) baselines completes zero full propose attempts
--- the first generation call itself does not finish within budget for most cells, so the 2k point
measures whether a truncated partial attempt happened to already contain a complete proof, not whether
the model got a fair single try. Attempts-per-cell only approach one full independent sample by $B$=8k
(mean propose count 0.82--0.87 across the four baselines) and keep climbing well past $N$=1 at every
larger budget we report (mean propose count at 128k: 1.94/2.22 miniF2F, 4.71/5.85 ProofNet\#, Goedel/
DeepSeek respectively; full table in \texttt{results/phase0/ATTEMPTS\_PER\_BUDGET\_TABLE.md}). Any
claim that maps this paper's budget axis onto a literature pass@$N$ number should use this table, not
assume $B$$\leftrightarrow$$N$ parity.

\section{The pass@\emph{B} curve: budget saturation and the in-distribution / OOD asymmetry}
\label{sec:passb}

\begin{table}[t]
  \centering
  \caption{pass@$B$, mean $\pm$ seed-std over 3 seeds. Same curve shape on both models: miniF2F
  saturates by 128k; ProofNet\# is still climbing from a low base.}
  \label{tab:passb}
  \begin{tabular}{lcccc}
    \toprule
    Budget & miniF2F $\cdot$ Goedel & miniF2F $\cdot$ DeepSeek & ProofNet\# $\cdot$ Goedel & ProofNet\# $\cdot$ DeepSeek \\
    \midrule
    2k   & 29.6 $\pm$ 3.3 & 27.9 $\pm$ 2.1 & 4.8 $\pm$ 1.4  & 5.4 $\pm$ 0.5 \\
    8k   & 60.1 $\pm$ 1.9 & 57.9 $\pm$ 1.7 & 9.3 $\pm$ 1.1  & 13.1 $\pm$ 0.8 \\
    32k  & 69.5 $\pm$ 0.6 & 67.1 $\pm$ 0.9 & 12.0 $\pm$ 0.6 & 18.3 $\pm$ 1.6 \\
    128k & 74.9 $\pm$ 0.9 & 72.0 $\pm$ 0.5 & 14.3 $\pm$ 0.8 & 22.2 $\pm$ 1.7 \\
    \bottomrule
  \end{tabular}
\end{table}

\begin{figure}[t]
  \centering
  \todo{F1: 2$\times$2 pass@B grid (miniF2F/ProofNet\# $\times$ Goedel/DeepSeek), generate from
  \texttt{results/*/metrics.json} via \texttt{scripts/analyze\_*.py}.}
  \caption{pass@$B$ curves, all four (model, benchmark) cells.}
  \label{fig:passb}
\end{figure}

Table~\ref{tab:passb} replicates on both models: miniF2F rises steeply from 2k to 8k ($\sim$+30pp) then
saturates to a $\sim$72--75\% ceiling by 128k; ProofNet\# is 3--5$\times$ harder at every budget and far
flatter, still climbing at 128k from a low base. This budget-saturation asymmetry --- in-distribution
saturates, out-of-distribution stays budget-hungry --- is not a single-model artifact.

A second, cross-model dichotomy is visible in the same table: Goedel slightly edges DeepSeek
in-distribution ($-2$ to $-3$pp on miniF2F), but DeepSeek clearly beats Goedel on the harder OOD
ProofNet\# at every budget, with the gap \emph{widening} with budget ($+8$pp at 128k). We rule out a
coverage artifact directly: both provers attempt the identical canonical statement sets under both Lean
pins (244/244 miniF2F, 186/186 ProofNet\# each compile on both pins), so the gap is not a
port/coverage difference. Tracing individual attempts (Section~\ref{sec:mechanism}) shows the OOD
advantage is deeper within-approach execution, not more approach diversity: every DeepSeek-only win is
on an opening Goedel also tried but could not close, and on shared misses DeepSeek diversifies
\emph{less} (3.08 vs 4.07 distinct openings) yet elaborates \emph{deeper} (median deepest step 29 vs
24). The lever that separates these two models is the model itself --- its training distribution and
recipe --- not search-time scaffolding, which is exactly the question the rest of this paper
interrogates directly.

\section{Agentic scaffolding is null (Phase 1)}
\label{sec:scaffolding}

We ablate six scaffolding components one-factor-at-a-time against a no-frills whole-proof + refinement
baseline (max\_iters 4, alloc\_split 0.5, all Phase 1 components off): BM25 premise retrieval ($k$=8),
a failure-memory buffer, an LLM reviewer over Lean-failed attempts, tactic-skeleton hints, and two
budget-allocation variants.

\textbf{miniF2F: no component clears the noise bar.} Retrieval looked like a mover in an initial pass
(+3.4pp@8k) but did not replicate under paired-flip inspection (+2.2/+0.7/$-0.8$pp at 2k/8k/32k; gains
$\approx$ losses at every budget --- BM25 context perturbs sampling without injecting usable premises).

\begin{table}[t]
  \centering
  \caption{ProofNet\# paired flips vs.\ the no-frills baseline. On the harder benchmark two components
  are directional, and both hurt.}
  \label{tab:flips}
  \begin{tabular}{lrrrl}
    \toprule
    Component & Gains & Losses & Net & Verdict \\
    \midrule
    reviewer            & 13 & 10 & $+3$  & noise-like \\
    budget\_alloc\_\_2   &  9 & 11 & $-2$  & noise-like \\
    memory               &  8 & 10 & $-2$  & noise-like \\
    tactic\_skeletons    &  9 & 12 & $-3$  & noise-like \\
    budget\_alloc\_\_0   &  6 & 25 & $-19$ & \textbf{directional (hurts)} \\
    retrieval             &  6 & 42 & $-36$ & \textbf{directional (hurts)} \\
    \bottomrule
  \end{tabular}
\end{table}

Table~\ref{tab:flips} sharpens the null on ProofNet\#: 4/6 components remain noise-like, but two are
directional --- and both harmful. Front-loaded budget allocation (no escalation) and BM25 retrieval (42
baseline solves lost for 6 gained) actively hurt out-of-distribution: bad context poisons the prompt
precisely when the model is already out of its depth. \textbf{No Phase~1 component is a net-positive
lever}, and the failure mode gets worse, not better, on the harder benchmark.

\section{Mechanism: why scaffolding cannot help, tested causally}
\label{sec:mechanism}

Section~\ref{sec:scaffolding} shows \emph{what}; this section explains \emph{why}, by mining the
on-disk attempt traces and then testing the resulting explanation with a pre-registered intervention,
not just correlation. Every finding below holds on both provers $\times$ both benchmarks.

\begin{itemize}
  \item[\textbf{F1}] \textbf{Diversity collapse.} On never-solved problems, the model tries only
    $\sim$1.8 distinct opening tactics across 18--33 attempts --- it loops on $\sim$2 approaches
    regardless of budget.
  \item[\textbf{F2}] \textbf{Failure taxonomy.} By our classifier's \emph{terminal} label (the single
    most severe failure mode per cell, across its full attempt history), unsolved attempts fail at deep
    reasoning (94--100\% \texttt{reasoning\_deep}) and almost never terminate on a missing lemma
    ($\sim$0--1\% \texttt{knowledge\_hallucinated\_lemma}). Because the terminal label reports only the
    last, most severe failure, it undercounts non-terminal premise/knowledge errors encountered earlier
    in a cell's history: on ProofNet\#, an unknown-identifier error is \emph{encountered} in 51.9\% of
    attempts even though it is almost never the terminal cause. Retrieval (as tested here, BM25 over a
    fixed premise corpus) still fails to help net of its context-perturbation cost (F2$'$,
    Section~\ref{sec:scaffolding}) --- but that failure is now better read as retrieval not solving a
    problem it does encounter, not as a distribution with nothing to retrieve; whether a stronger,
    neural retriever would fare differently is untested (Section~\ref{sec:related}).
  \item[\textbf{F3}] \textbf{Capability floor.} Unsolved attempts reach a median deepest step of
    26--53 and almost never stall at step 1: the model makes real progress into a proof but cannot
    \emph{close} it.
  \item[\textbf{F4}] \textbf{No easy subfield.} ProofNet\# solve rates are 21--33\% across every
    subfield --- there is no soft target a smarter router could farm.
  \item[\textbf{F5}] \textbf{Late solves re-use the same approach (correlational).} When a hard problem
    is finally solved late (after $\geq$3 attempts), it comes from re-sampling the \emph{same} approach,
    not a newly explored one (0\% new opening). This predicts that forcing approach diversity will not
    help, but is survivor-conditioned and cannot establish causation on its own.
\end{itemize}

\subsection{F6 / Step C: the interventional test (centerpiece)}
\label{sec:stepc}

We built an approach-conditioned diversity-injection component --- lists the distinct openings already
tried, asks explicitly for a fundamentally different approach --- and ran it on each model$\times$
benchmark's \emph{trapped core} (problems unsolved by every seed at the 128k baseline) at 8k/32k, 3
seeds, all four cells. We pre-registered the prediction: diversity rises, solves stay flat.

\begin{figure}[t]
  \centering
  \todo{F4: Step C before/after diversity + solve panel, all four cells, generate from
  \texttt{results/phase6/} STAGE\_C / diversity-injection artifacts.}
  \caption{Step C manipulation check (diversity) and outcome (solves), before/after, all four cells.}
  \label{fig:stepc}
\end{figure}

\textbf{The manipulation check passed}: budget-matched (@32k, attempt counts held fixed), distinct
openings per attempt rose 42--70\% ($\sim$1.0--1.3 $\to$ $\sim$1.7--1.8) --- the intervention genuinely
fired, not just nominally. \textbf{Solves stayed null}: trapped pass@8k/32k $\approx$ 0 everywhere
(Goedel miniF2F 0.6/1.2\%, Goedel ProofNet\# 0/0.7\%, DeepSeek both 0/0\%) --- exactly 5 genuine
verified flips across all four arms (0 on DeepSeek), within seed-std. \textbf{Output quality mildly
degraded}: last-attempt failure modes shift from \texttt{reasoning\_deep} (38--88\%) to majority
\texttt{formalization\_syntax} (62--74\%) plus \texttt{loophole\_sorry} (23--36\%); soundness-relevant
rates rise ($\sim$3$\times$ loophole, $\sim$1.5--2$\times$ syntax). Pushed for novelty, the model leaves
its competent manifold --- but every degraded attempt is caught by the verifier (0 false solves), so
result soundness is intact and the intervention is measurably counterproductive, not merely neutral.

This closes F5's correlational gap with a causal experiment on two independent provers:
\textbf{approach discovery is not the bottleneck; within-approach execution depth (the F2/F3 reasoning
floor) is.} No downstream lever that only changes \emph{which} approach is tried, without changing the
model's capacity to execute a chosen approach to completion, should be expected to move solve rate ---
which sets up the exhaustion sweep in Section~\ref{sec:exhaustion}.

\section{Exhaustion sweep: five further candidate levers, all null}
\label{sec:exhaustion}

Having isolated the mechanism causally, we exhaust the remaining candidate levers a reviewer might
propose, compressed here (full detail in the phase result docs cited inline).

\textbf{Hammer/SMT closing tactics (Phase 3).} A portfolio of hammer/SMT-style closing tactics against
the trapped core: 0/119 newly solved. Rules out ``the model just needs a stronger closing-tactic
library'' as cheaply as Section~\ref{sec:scaffolding} ruled out scaffolding.

\textbf{More of the same budget (Phase 5).} Extending already-trapped cells past 128k tokens (more of
the \emph{same} budget on the \emph{same} problems): $\sim$0 new solves (Goedel 1, DeepSeek 0). Budget
helps only when reallocated \emph{across} different problems (Section~\ref{sec:allocation}), not by
adding more of it to a problem already known to be hard --- a fourth, independent confirmation of the
execution floor.

\textbf{Closing-likelihood SFT (Phase 6, Stage A/B).} Fine-tuning to maximize conditional likelihood on
verified proof closings is a two-model null, but a diagnostically useful one: near-zero teacher-forced
loss on the closing, yet the same model still fails when it has to \emph{generate} its own path there
autoregressively --- an exposure-bias signature. This reframes the open question from ``does the model
know how to close a proof'' (yes) to ``why does it drift away from a closeable state during its own
free-running generation'' (open).

\textbf{Lightweight post-hoc RL (Phase 6, Stage C).} A GRPO probe (LoRA $r$=16, DeepSeek-Prover-V2-7B,
80 steps, direct outcome reward against the Lean verifier) is a clean capacity-ceiling null: held-out
pass@1 moved $0.586 \to 0.570$ ($\Delta = -1.6$pp), training reward flat across all 80 steps, with
manipulation checks confirming no reward-hacking or diversity collapse. RL, at LoRA scale, applied
post-hoc, does not move the floor.

\textbf{Verified-state re-grounding (Phase 7).} Tests the Stage~B exposure-bias hypothesis directly and
training-independently: force continuation from a \emph{verified} intermediate Lean state rather than
free-running generation. Both re-grounding modes tested, both models: null. Re-grounding alone does not
unlock the trapped core.

\textbf{Full-pipeline, lab-scale RL (Phase 8).} The strongest remaining candidate: not a lightweight
post-hoc probe but a lab's own full Base$\to$SFT$\to$RL pipeline, tested on two independently matched
training lineages (DeepSeek-Prover-V1.5 and Leanabell-Prover). This phase surfaced and fixed four real,
independently confirmed harness bugs before its result could be trusted (Section~\ref{sec:methodology});
after correction and a harness-sanity control (37/37 and 40/40 historically-solved proofs still verify
under the patched backend), the corrected result is \textbf{0.0 $\pm$ 0.0 pass@$B$ at every training
stage, every budget up to 32000, both benchmarks, both lineages} --- the originally reported
Base\,$<$\,SFT\,$<$\,RL pattern was entirely a wiring-bug artifact and does not survive correction.

\section{The constructive counterpoint: compute-optimal budget allocation}
\label{sec:allocation}

The one lever that moved anything in this project is not a change to any single problem's search or
training at all, but a \emph{policy} for how a fixed total token budget is split across a batch of
\emph{different} problems.

\begin{figure}[t]
  \centering
  \todo{F5: allocation result --- efficiency frontier (realizable vs.\ uniform vs.\ oracle),
  \texttt{results/phase4/frontier\_*.png}.}
  \caption{Compute-optimal allocation: realizable policy vs.\ uniform vs.\ oracle ceiling,
  Goedel$\times$ProofNet\#.}
  \label{fig:allocation}
\end{figure}

On ProofNet\# ($\sim$80\% of cells never solve, so uniform allocation pours full budget into a
provably-trapped core), a \emph{realizable} during-run policy --- abandon cells a logistic predictor
flags as trapped at a decision checkpoint $c^*$, informed by the F2/F3 mechanism (elapsed spend, proof-
depth plateau) --- saves substantial compute at near-equal accuracy on Goedel: \textbf{+26\% $\pm$ 7\%}
saved at 90\% of full-budget accuracy, robust across all three independent seeds (+25/+18/+34). The same
policy on DeepSeek is weak and fragile by comparison ($-13\% \pm 28\%$, driven by a seed-2 collapse to
$-51\%$) --- \textbf{this positive result is one-model-robust, not yet a confirmed general property}.
Our own post-hoc mechanism analysis (\texttt{ALLOCATION\_MECHANISM.md}) finds this fragility is better
explained as a thin-population per-seed sampling artifact than a structural deficit in DeepSeek's
population of exploitable cells, which motivates an 8-seed power-up \todo{report resolved DeepSeek
result once WS1.1 GPU data lands} as the natural next confirmatory step. We report this result honestly
scoped: the win is at a fractional-accuracy operating point (at 100\% accuracy the policy must keep
almost every cell, since the costliest winnable cells are indistinguishable from trapped ones at the
decision checkpoint, so the saving is $\sim$0 there by construction, not by tuning failure), and
efficiency rather than raw fixed-compute accuracy is the robust axis.

\section{Methodology: auditing the verifier, and symmetric skepticism}
\label{sec:methodology}

Every number in this paper depends on the Lean verifier being sound, so we treat it as a first-class
audited artifact rather than an assumed-correct oracle, across two independent audit rounds.

\textbf{Round 1 (2026-06-14).} Found two verifier holes producing false-positive solves on hard
problems: a truncated generation emitting a bare \texttt{def} preamble that Lean compiles with no goal
(scored solved), and a wedged/cross-talked REPL returning an empty no-error response (scored success).
Both fixed; blast radius measured directly rather than assumed: miniF2F Phase~0 had 0 false positives,
Phase~1 had 6/3124 (0.2\%); only ProofNet\# was materially corrupted (constant full-budget truncation),
and was fully re-run on the fixed verifier.

\textbf{Round 2 (independent audit, 2026-07-10/16).} A separate audit pass, run before any further
experiments were queued, found a further soundness gate bug: the \texttt{no\_goal} check compared the
model's raw extracted completion against a regex, rather than the backend's actually-assembled/compiled
source --- structurally always-failing for continuation-style completions regardless of correctness.
Confirmed, via a taint audit and three independent from-scratch re-derivations that matched the
committed pass@$B$ curve, oracle ceiling, and flip table exactly, to be Phase-8-only and to have zero
effect on any number reported in Sections~\ref{sec:passb}--\ref{sec:allocation}
(\texttt{whole\_proof}, the template family both headline models use, is a no-op case for this gate).
A further, materially small (13/1212 cells, 2.0\% of trapped problems, both models) retroactive
correction was found on re-verifying trapped-core failures under a Lean elaboration-heartbeat setting
introduced after those results were generated; every flipped cell's proof was already present in the
original generation, so the fix only ever widens what counts as solved and the reported floor is, if
anything, a small upward correction, not an overturned result. \todo{fold the 13 corrected cells back
into the affected pass@B curves once flowed --- arithmetic-only, not a new experiment.} This 13/1212
figure is the \emph{final-scoring} blast radius, and understates the setting's live effect on the
refinement loop itself: on miniF2F, 17.8\% of refine steps (both models, 550/3090 Goedel and
718/4046 DeepSeek) were immediately preceded by a spurious Lean heartbeat timeout that the
model then had to react to as if it were a real proof error --- a live trajectory distortion, not
just a scoring artifact, that this paper has not attempted to quantify beyond the terminal-cell flip
count. \textbf{Methods warning for refinement-loop builders:} an elaboration-heartbeat limit that is
lenient at verification time but strict (or absent) during the same run's intermediate refine calls
silently injects false negative feedback into the loop the model is conditioning on --- worth an
explicit audit before trusting any refinement-based agent's feedback signal, independent of whether
the final scoring bug it causes turns out to be small.

\textbf{Lessons carried forward, with sharpest first:} (1)~paired flips beat mean-deltas at $n{=}3$
seeds --- the flip table is what exposed retrieval as churn rather than signal. (2)~A second, harder,
OOD benchmark earns its cost: it corroborated the miniF2F null and surfaced components (retrieval,
front-loaded allocation) that actively hurt, invisible on the easier benchmark. (3)~A second,
independently trained model earns its cost: it turned a single-lab result into a mechanism that
generalizes and surfaced the in-distribution/OOD dichotomy. (4)~Distinguish the manipulation check from
the outcome: Step~C's null is only interpretable because we first confirmed the intervention fired
(budget-matched distinct-opening counts), not raw attempt-confounded diversity. (5)~\textbf{Apply the
same base-rate skepticism to a surprising null as to a surprising positive.} Phase~8's
0.0\%-everywhere floor looked like a genuine, clean Base$<$SFT$<$RL finding until it was checked against
the models' own published pass rates and a byte-level prompt diff --- the discipline that would catch an
implausibly high number should catch an implausibly low or implausibly perfect one with equal rigor. A
cheap harness-sanity control against an unaffected reference model (Section~\ref{sec:exhaustion},
Phase~8) should run \emph{before} trusting a dramatic null, not only before trusting a dramatic positive.
(6)~A newly introduced model configuration exercises code paths existing tests never covered: all three
Phase~8 structural bugs were latent in code that had shipped, and presumably passed review, for weeks,
because no test exercised a template family other than \texttt{whole\_proof} until Phase~8 needed one at
real scale --- add the new case as a failing test against real data \emph{first}.

\section{Discussion, limitations, and anticipated objections}
\label{sec:discussion}

\textbf{``Your harness is broken.''} Section~\ref{sec:methodology} is not an appendix afterthought: the
Phase~8 bug-hunt, the 37/37 and 40/40 harness-sanity controls, and two independent verifier audits are
reported in the main text because a floor result is only as credible as the verifier that produced it.

\textbf{``Only two models.''} We additionally test four training lineages, including two matched
Base/SFT/RL sequences (Section~\ref{sec:exhaustion}, Phase~8) --- precisely to test
\emph{training-stage} effects, not merely to add datapoints to the pass@$B$ curves.

\textbf{``Scaffolding was implemented weakly.''} The Section~\ref{sec:stepc} manipulation check is the
direct answer: the diversity intervention demonstrably fired (42--70\% more distinct openings,
budget-matched) and still did not move solves.

\textbf{``miniF2F is saturated/contaminated.''} This is exactly why ProofNet\# carries the
out-of-distribution claims in this paper; the saturation/still-hungry dichotomy between the two
benchmarks is reported as a finding, not hidden as a limitation.

\textbf{``Negative results aren't contributions.''} We position this paper directly against an active
test-time-compute and agentic-scaffolding literature that is shipping the same components
(Section~\ref{sec:related}) this paper shows are null or actively harmful for whole-proof theorem
proving specifically, with a causal (not correlational) mechanism account and two independently
matched training-lineage negatives to back it.

\textbf{``The mathlib pin is stale, so `unknown identifier' errors are really a harness artifact.''}
Our verifier is pinned to a fixed mathlib4 snapshot (Section~\ref{sec:methodology}), and the prover
does hit unresolved identifiers at a non-trivial rate on ProofNet\# (51.9\% of attempts encounter one,
Section~\ref{sec:mechanism}), which invites this objection. We checked it directly: of 20 uniformly
sampled unresolved identifier names from ProofNet\# baseline attempts, only 3/20 (15\%) have a
same-named declaration anywhere in a fresh clone of current upstream mathlib4, and even those three
are generic, commonly-reused lemma-name suffixes rather than confirmed matches in the right namespace.
The remaining 17/20, including the single most frequent unresolved name in the corpus
(\texttt{IsGroupHomomorphism}, 27 occurrences), have no match at all --- consistent with the model
targeting an API surface (unbundled \texttt{Is*Hom} predicates) that mathlib4 has since replaced
outright, not merely renamed. This sample leans toward genuine model/API mismatch over pin staleness,
though it does not by itself distinguish training-data mathlib version from harness pin version as the
source of that mismatch.

\textbf{Open question.} The exposure-bias signature (Section~\ref{sec:exhaustion}, Phase~6 Stage~B) ---
near-zero teacher-forced loss on a verified closing, yet autoregressive failure to reach it --- is this
project's most interesting unresolved thread and is analysis-, not training-, shaped; a drift
characterization is ongoing background work (\texttt{results/phase9/DRIFT.md}) and, if a crisp
signature emerges, a natural follow-up.

\section{Conclusion}

Across eight phases, two independently trained provers, two benchmarks with an in-distribution/OOD
asymmetry, and four training lineages including two matched Base$\to$SFT$\to$RL sequences, we find a
consistent execution floor: agentic scaffolding, additional search, additional same-problem budget,
closing-likelihood SFT, lightweight post-hoc RL, and full-pipeline lab-scale RL all fail to move solve
rate, and a pre-registered causal intervention confirms directly that the bottleneck is within-approach
execution depth, not approach discovery. The one lever that reliably moves outcomes at fixed compute is
allocation policy across a batch of different problems --- reported honestly as one-model-robust, with
its own mechanism account and a stated path to confirmation. We hope this serves both as a cautionary
result for test-time-compute claims in formal mathematics specifically, and as a methodological template
--- verifier auditing, causal not correlational mechanism tests, and symmetric skepticism of nulls and
positives alike --- for negative-result reporting in this area more broadly.

\bibliographystyle{plainnat}
\bibliography{refs}

\end{document}
